\documentclass{beamer}
\usetheme{CambridgeUS}

\usepackage{tikz}
\usetikzlibrary{calc,positioning}

\title{Arbitrage-Free Pricing Theory}
\author{Matteo Sani}
\begin{document}
	\begin{frame}[plain]
		\maketitle
	\end{frame}

\AtBeginSection{
\begin{frame}{Outline}
	\tableofcontents[currentsection]
\end{frame}
}
\section{Fundamental Definitions}
\subsection{Portfolio, Numeraire, Arbitrage\ldots}
\begin{frame}{Portfolio}
	\begin{block}{Portfolio Definition}
		A \textcolor{red}{portfolio} is a vector $\mathbf{\theta}\in \mathbb{R}^K$ whose $j$ components represent the number of shares of asset $A_j$. It's value is
		\begin{equation}
			V_t(\mathbf{\theta}, \omega)=\sum_{j=1}^K\theta_jS^j_t(\omega)
		\end{equation} 
		where $S_t^j$ is the value of $j$th asset, and $\omega$ a market situation. A portfolio is \textcolor{red}{self-financing} if its value changes only due to variations of the asset prices.
	\end{block}
\begin{block}{Numeraire Definition}
A \textcolor{red}{numeraire} is any positive non-dividend-paying asset. It is a reference asset chosen to normalize all other asset prices to it. Having a numeraire allows for the comparison of the value of goods against one another.
\end{block}
\end{frame}

\begin{frame}{Arbitrage}
	\begin{block}{Arbitrage Definition}
		An \textcolor{red}{arbitrage} is a self financing portfolio that \emph{makes money from nothing}, formally a portfolio $\mathbf{\theta}$ such that
		\begin{equation}
				V_0(\mathbf{\theta}, \omega)\le 0 \text{ and } \mathcal{P}\{V_t(\mathbf{\theta}, \omega) > 0\} > 0
		\end{equation}
	\end{block}
	Informally, an arbitrage is a way to make a guaranteed profit from nothing, by short-selling certain assets at time $t = 0$, using the proceeds to buy other assets, and then settling accounts at time $t$.
\end{frame}

\subsection{Fundamental Theorems of Arbitrage Pricing}
\begin{frame}{Fundamental Theorem of Arbitrage Pricing (I)}
	\begin{block}{Risk-Neutral Measure Definition}
		A probability distribution (measure) $\mathbb{N}$ on the set $\Omega$ of possible market scenarios is said to be a \textcolor{red}{risk-neutral measure} if, for every asset $A$, and numeraire $N$ the price of $A$ at time $t=0$ is the \emph{discounted expectation}, under $\mathbb{N}$, of the price at time $t$, that is
		\begin{equation}
			\frac{S_0}{N_0} = \mathbb{E}^{\mathbb{N}}\left[\frac{S_t}{N_t}| \mathcal{F}_t\right]
			\label{eq:risk_neutral_pricing}
		\end{equation}
		The discounted value process is a martingale.
	\end{block}
	\begin{block}{Fundamental Theorem of Arbitrage Pricing I}
		There exists a risk-neutral measure if and only if arbitrages do not exist.
	\end{block}
\end{frame}

\begin{frame}{Hedging}
	\begin{itemize}
		\item A portfolio $\mathbf{\theta}$ in the assets $\mathbf{A}$ is a \textcolor{red}{replicating portfolio} for the asset $B$ if
		\begin{equation}
			S_t^{B}(\omega_i) = \sum_{j=1}^K \theta_j S_t^j(\omega_i)\quad\forall i=1,2,\ldots,N
		\end{equation}
		\item In particular if the market is \emph{arbitrage-free} then the relation holds for all $t$.
		\item The importance of replicating portfolios is that they enable financial institutions that sell
		asset $B$ (e.g. a call options) to \textcolor{red}{hedge}: for each sold share of asset $B$, buy $\theta_j$ shares
		of asset $A_j$ and hold them to time $t + 1$. Then at time $t + 1$, 
		\begin{equation*}
			\text{net gain }= \text{ net loss } = 0
		\end{equation*}
	\end{itemize}
\end{frame}

\begin{frame}{Fundamental Theorem of Arbitrage Pricing (II)}
	\begin{itemize}
		\item In some circumstances, an arbitrage-free market may admit more than one risk-neutral measure, i.e. \textcolor{red}{incomplete markets}.
		\item By contrast, a \textcolor{red}{complete market} is one that has a unique risk-neutral measure.
	\end{itemize}
	\begin{block}{Completeness Theorem Definition}
		Let $\mathcal{M}$ be an arbitrage-free market with a risk-less asset. If for every derivative security there is a replicating portfolio in the assets $A_j$ then the market $\mathcal{M}$ is complete. Conversely, if the market $\mathcal{M}$ is complete, and if the unique risk-neutral measure $\mathbb{P}$ gives positive probability to every market scenario $\omega$, then for every
		derivative security there is a replicating portfolio in the assets $A_j$.
	\end{block}
\end{frame}

\begin{frame}{Summary of Basic Definitions}
	\begin{itemize}
		\item The market is free of arbitrage if (and only if) there exists an \textcolor{red}{equivalent martingale measure} (i.e. a risk-neutral measure).
		\item The market is complete if and only if the martingale measure is unique.
		\item In an arbitrage-free market the price of any derivative is uniquely given, either by the value of the associated replicating strategy, or by the risk-neutral expectation of the discounted payoff under the risk-neutral measure. 
	\end{itemize}
\end{frame}

\subsection{Money Market Account}
\begin{frame}{Money Market Account}
	\begin{itemize}
		\item A \textcolor{red}{money market account} represent a risk-less investment, where profit is accrued continuously at the risk-free rate, and its value is denoted by $B(t)$.
		\item We assume $B(0)=1$ and that 
		\begin{equation}
			dB(t) = r(t)B(t)dt
		\end{equation}
		\item This evolution process has the solution 
		\begin{equation}
			B(t) = \exp\left(\int_0^t r_s ds\right)
		\end{equation}
		where $r_t$ is referred to as \textcolor{red}{instantaneous spot rate} or  \textcolor{red}{short rate}.
	\end{itemize}
\end{frame}

\begin{frame}{Money Market Account}
	\begin{itemize}
		\item In case $r_t$ is deterministic, from the definition of money market account it follows that 
		\begin{equation*}
			V(0) = A \implies V(t) = A\cdot B(t)
		\end{equation*}
		\item If we want at time $T$ exactly 1 unit of currency
		\begin{equation*}
			AB(T) = 1 \implies AB(t) = \frac{B(t)}{B(T)} 
		\end{equation*}
		hence $\frac{B(t)}{B(T)}$ is the value of one unit of currency payable at time $T$ seen from $t$.
		
		%%		\item We now define the abstract quantity $r(t)$, the \textbf{short rate}, as
		%%		\begin{equation}
			%%			r(t) = \lim_{T\rightarrow t^+} L(t,T) \simeq L(t, t+\epsilon)
			%%		\end{equation}\quad with $\epsilon$ small
	\end{itemize}
\end{frame}

\subsection{Stochastic Discount Factor and Zero Coupon Bond}
\begin{frame}{Stochastic Discount Factor}
	\begin{itemize}
		\item The \textcolor{red}{(stochastic) discount factor} $D(t, T)$ is the amount at time $t$ that is \emph{equivalent} to one unit of currency payable at time $T$ and is given by
		\begin{equation}
			D(t, T) = \frac{B(t)}{B(T)} = e^{-\int_t^T r(s) ds}
		\end{equation}
		\item Can you guess which are its dimension ?	
		\item In many pricing application (e.g. Black-Scholes formula) $r$ is assumed to be a deterministic function of time, and so are $B(t)$ and $D(t,T)$.
		\begin{itemize}
			\item This is motivated by the small influence interest rate variations have on equity prices.
		\end{itemize}
		\item When dealing with interest rate products, $r$ becomes the main actor, so the deterministic assumption must be dropped.
	\end{itemize}	
\end{frame}

\begin{frame}{Money Market Account as Numeraire}
	If we choose the money market account $B(t)$ as numeraire, then Eq.~\ref{eq:risk_neutral_pricing} becomes
	\begin{equation*}
		\begin{aligned}
			S_t &= \mathbb{E}^{\mathbb{Q}_B}\left[\frac{B_t}{B_T}S_T|\mathcal{F}_t\right]\\
			&=\mathbb{E}^{\mathbb{Q}_B}\left[e^{-\int_t^T r_s ds}S_T|\mathcal{F}_t\right]
		\end{aligned}
	\end{equation*}
	where $\mathbb{Q}_B$ is the risk-neutral measure that makes $\frac{S_t}{B_t}$ a martingale. 
	
	When interest rates are deterministic, the exponential can be brought out of the expectation
	\begin{equation*}
		S_t = e^{-\int_t^T r_s ds} \mathbb{E}^{\mathbb{Q}_B}\left[S_T|\mathcal{F}_t\right]
	\end{equation*}
	And finally when rates are constant (e.g. Black-Scholes or Heston models)
	\begin{equation*}
		S_t = e^{-r(T-t)}\mathbb{E}^{\mathbb{Q}_B}\left[S_T|\mathcal{F}_t\right]
	\end{equation*}
\end{frame}

\begin{frame}{Zero Coupon Bond}
	\begin{itemize}
		\item \textcolor{red}{Zero Coupon Bond} (ZCB) is a contract that pays one unit of money at time $T$. Its price at time $t$ is denoted by $P(t,T)$, and by definition $P(T,T) = 1$.
		\item What is the relation between $P(t,T)$ and $D(t,T)$ ? \\
		\renewcommand{\arraystretch}{1.2}
		\footnotesize{\tiny {\tiny }}{
			\begin{table}[bt]
				\begin{tabular}{|c|c|} \hline
					$D(t, T)$ & $P(t, T)$ \\ \hline
					equivalent amount of money & value of a contract \\ \hline
					\multicolumn{2}{|c|}{deterministic rates} \\ \hline
					\multicolumn{2}{|c|}{D(t,T)=P(t,T)} \\ \hline 
					\multicolumn{2}{|c|}{stochastic rates} \\ \hline
					random quantity at time $t$ &
					\renewcommand{\arraystretch}{1.0} 
					\begin{tabular}{@{}c@{}}
						being the value of a contract\\ with payoff at time $T$ must be known in $t$ \\
					\end{tabular} \\ \hline
				\end{tabular}
			\end{table}
		}
		\item We will see that in absence of arbitrage
		\begin{equation*}
			P(t, T) = \mathbb{E}^Q[D(t, T)|\mathcal{F}_t]
		\end{equation*}	
	\end{itemize}
\end{frame}

\begin{frame}{Intermezzo}
	\begin{itemize}
		\item Time to maturity: $T-t$
		\item Year fraction, day-time convention. Denote by $\tau(t, T)$ the chosen time measure between $t$ and $T$. In theory, you can think of $\tau(t, T)=T-t$. In practice, among the plethora of convention, the following are worth mentioning:
		\begin{itemize}
			\item Actual/365
			\item Actual/360
			\item 30/360
		\end{itemize}
		\item \textcolor{red}{Go and find their definitions: they are also embedded in Excel as financial functions.}
	\end{itemize}
	\begin{tikzpicture}[remember picture,overlay]
	\node[xshift=5cm,yshift=-3.7cm] (image) at (current page.center) {\includegraphics[width=20px]{python_logo}};
	\node[align = center, yshift=1.45cm, below=of image] {\tiny{\href{shorturl.at/chzT2}{shorturl.at/chzT2}}};
	\end{tikzpicture}
\end{frame}

\subsection{Spot Interest Rate}
\begin{frame}{Compounding}
	\begin{itemize}
		\item Simply-compounded \textcolor{red}{spot interest rate}
		\begin{equation}
			L(t,T)=\frac{(1-P(t,T))}{\tau(t,T)P(t,T)}		
		\end{equation}
		\item The \textcolor{red}{yield curve} at time $t$ is the graph of:
		\begin{equation*}
			T\rightarrowtail L(t, T)
		\end{equation*}
		\item In the market (and in the \emph{books}) these are the so called LIBOR (or EURIBOR) rates and are typically compounded with the actual/360 convention. They are the main rates underlying interest rate derivatives.
	\end{itemize}
\end{frame}

\begin{frame}{Compounding}
	In the following we recall some useful definition \textbf{you should be always familiar with}.
	\begin{itemize}
		\item Annually-compounded spot interest rate
		\begin{equation}
			Y(t, T)= \frac{1}{P(t, T)^{\frac{1}{\tau(t,T)}}}-1
		\end{equation}
		\item $k$-times-per-year compounded spot interest rate
		\begin{equation}
			Y^k(t, T)= \frac{k}{P(t, T)^{\frac{1}{k\tau(t,T)}}}-1
		\end{equation}
		\item When $k\rightarrow\infty$, we get the continuously compounded rate
		\begin{equation}
			R(t,T)=-\frac{\log P(t,T)}{\tau(t,T)} \implies P(t,T)=e^{-R(t,T)\tau(t,T)}
		\end{equation}
	\end{itemize}
\end{frame}

%\begin{frame}{Fiddling with Rates (2/2)}
%	MIGLIORARE
%	\begin{itemize}
	%	\item Does it ring a bell ?
	%	\item We can also write
	%	\begin{equation} 
		%	f(t, T)=-\frac{\partial \log P(t, T)}{\partial T} = -\frac{\partial\log \mathbb{E}^Q[e^{-\int_t^T r(s) ds} | \mathcal{F}_t]}{\partial T}
		%		\end{equation}
	%	\item So what is the difference between $r(t)$ and $f(t, T)$ ?
	%	\end{itemize}
%\end{frame}

\section{Forward Rate Agreement}
\begin{frame}{Forward Rate Agreement: Definition}
	\begin{itemize}	
		\item A \textcolor{red}{Forward Rate Agreement} (FRA) is a contract involving three time instants: %the current time $t$, the fixing time $T>t$, and the maturity time $S>T$.
		\begin{equation*}
			\underbrace{t}_{\text{current time}} \leq \underbrace{T}_{\text{fixing time}} \leq\underbrace{S}_{\text{maturity}}
		\end{equation*}
		\item The FRA payout consists of an exchange of interest rate flows calculated for the time period $\tau=S-T$. At the maturity $S$, a fixed payment based on a fixed rate $K$ is exchanged against a floating payment based on the spot rate $L(T, S)$ whose value is known only in $T$.
		\item Basically, this contract allows one to lock-in the interest rate between times $T$ and $S$ at a desired value $K$. Interest rate flows are calculated using the simple compounding law.
	\end{itemize}
\end{frame}

\begin{frame}{FRA: Formalization of the Contract}
	\begin{itemize}
		\item Formally, at time $S$ one receives $\tau(T, S)KN$ units of currency and pays the amount $\tau(T,S)L(T,S)N$, where $N$ is the contract nominal value.
		\item Thus, at time $S$, the (future and today unknown) payout of the contract is: 
		\begin{equation}
			N\tau(T,S)(K-L(T,S))
			\label{eq:fra_payoff}
		\end{equation}
		\item In order to assign a value at this contract we have to tackle two issues:
		\begin{itemize}
			\item how to estimate $L(T, S)$
			\item how to discount it from $S$ to $t$
		\end{itemize}.
		\item There are several ways to arrive at the final result: \textbf{no arbitrage is the common denominator.}
	\end{itemize}
\end{frame}

%\begin{frame}{FRA: Numerical Example}
%	A company enters into a FRA to receive 4\% on \$100M for a 3-month period starting in 3 years. Imagine that LIBOR is 4.5\% in 3 years.
%	
%	So $T=3y$, $S=3.25y$, $K=4\%$, and $L(T,S)=4.5\%$.
%	
%	The net cash-flow at time $S$ is
%	\begin{equation*}
%		100000000\times(0.04-0.045)\times0.25 = -\$125000
%	\end{equation*} 
%\end{frame}

\subsection{FRA Valuation}
\begin{frame}{FRA: the Standard Derivation}
	\begin{itemize}
		\item Following the approach reported in Brigo-Mercurio (2006) we can proceed as follows:
		\item Substitute in Eq.~\ref{eq:fra_payoff} $L(T,S)$ with its expression as a function of $P$ and get
		\begin{equation*}
			N[\tau K - \frac{1}{P(T, S)} + 1]
		\end{equation*}
		\item Interpret $A=\frac{1}{P(T,S)}$ as an amount of money held at time $S$. At $T$ it's worth 1, which in turn, at time $t$, equals $A=P(t,T)$.
		\item On the other hand $B=\tau K + 1$ in $S$ becomes at time $t$
		\begin{equation*}
			P(t,S)\tau K + P(t, S)
		\end{equation*}
		\item Collecting the terms we get
		\begin{equation}
			\mathbf{FRA}(t,T,S,\tau,N,K)=N[P(t,S)\tau K–P(t,T)+P(t,S)]
		\end{equation}
	\end{itemize}
\end{frame}

\subsection{FRA Valuation ($2^{nd}$ Approach)}
\begin{frame}{FRA Valuation: a Different Approach}
	Now we will try to estimate the value of a FRA using a simple replication approach:
	\begin{itemize}
		\item at time $t$: 
		\begin{equation*}
			\begin{cases}
				\text{lend}~P(t,T)\\
				\text{borrow}~P(t,S)(1+(S-T)K)
			\end{cases};
		\end{equation*}
		\item at time $T$: you receive 1 which you will invest at the (in $t$ unknown) EURIBOR rate $L(T,S)$;
		\item at time $S$: 
		\begin{equation*}
			\begin{cases}
				\text{receive}~(1+L(T,S)(S-T))\\
				\text{pay}~1 + (S-T)K
			\end{cases}.
		\end{equation*}
	\end{itemize}
\end{frame}

\begin{frame}{FRA Valuation: a Different Approach}
	Adding the cash-flows together and evaluating in $t$ we get:
	\begin{equation*}
		\begin{gathered}
			1+L(T,S)(S-T) - 1 - (S-T)K = (L(T,S)-K)(S-T) \\
			\mathbb{E}^Q[D(t, S)(L(T, S)-K)(S-T)]
		\end{gathered}
	\end{equation*}
	
	But its value must equal to the value of the replicating portfolio in $t$ (we are in a complete market\ldots). Hence
	\begin{equation}
		\begin{aligned}
			\mathbb{E}^Q[D(t,S)&(L(T,S)-K)(S-T)]=\\
			&[P(t,S)(S-T)K-P(t,T)+P(t,S)]
		\end{aligned}
	\end{equation}
\end{frame}

\subsection{Forward Rate Definition}
\begin{frame}{FRA: the Forward Rate as a Break-even Rate}
	\begin{block}{Forward Rate}
		The value of $K$ which makes the value of the FRA equal to zero is called \textcolor{red}{simply-compounded forward rate}
		\begin{equation}
			F(t;T,S):=\frac{1}{\tau(T,S)}\left[\frac{P(t,T)-P(t,S)}{P(t,S)}\right]
		\end{equation}
	\end{block}
	
	The forward rate can be interpreted as a rate observed in $t$ spanning the time period $S-T$.
	
	%Its value depends on no-arbitrage consideration.
	
\end{frame}

\begin{frame}{Forward Rate}
	\begin{itemize}
		\item We can then rewrite the value of the FRA in terms of the simply-compounded forward rate
		\begin{equation}
			\mathbf{FRA}(t,T,S,\tau,N,K)=NP(t,S)\tau(T,S)(K-F(t;T,S))
			\label{eq:fram_payoff_withF}
		\end{equation}
		(this formula will be used for the swap evaluation)
		\item It is just like we had replaced the LIBOR rate $L(T,S)$ in the payoff with the forward rate $F(t,T,S)$ and then taken the present value of the (deterministic) quantity.
		\item The forward rate can be interpreted as an \emph{estimate} of the future spot rate, which is random at time $t$ (based on the market conditions).
		\item We'll see later that indeed $F(t,T,S)$ is the expectation of $L(T,S)$ under a particular probability measure.
	\end{itemize}
\end{frame}

\begin{frame}{Instantaneous Forward Rate}
	\begin{itemize}
		\item Now we introduce the \textcolor{red}{instantaneous forward rate} $f(t, T)$.
		It is defined as 
		\begin{equation}
			\begin{aligned}
			f(t, T) &:= \lim_{S\rightarrow T^+} F(t,T,S) \\
			& = \lim_{\epsilon\rightarrow 0}  \frac{1}{\tau(T,T+\epsilon)}\frac{P(t,T)-P(t,T+\epsilon)}{P(t,T+\epsilon)} \\
			& = \lim_{\epsilon\rightarrow 0} - \frac{1}{P(t,T)} \frac{P(t,T+\epsilon)-P(t,T)}{\epsilon} \\
			& = -\frac{\partial \log P(t, T)}{\partial T}
			\end{aligned}
		\end{equation}
	\end{itemize}
\end{frame}

\begin{frame}{Instantaneous Forward Rate}
	\begin{itemize}
		\item From the previous equation we can derive
		\begin{equation}
			P(t, T) = e^{-\int_t^T f(t, s) ds}
		\end{equation}
		\item The instantaneous forward rate is the forward rate for a forward contract with an infinitesimal investment period after the settlement date.
		\item Notice that
		\begin{equation}
			r(t) = f(t,t)
		\end{equation}
	\end{itemize}
\end{frame}

%\begin{frame}{}
%\begin{block}{}
%This derivation seems very general but in reality it relies on some strong assumptions about the absence of counter-party liquidity credit risk. If that book were written today may be the authors would give a different version of the proof. The next derivation is a bit more explicit in the assumptions needed to get to the result.
%\end{block}
%\end{frame}
%
%\begin{frame}{Derivation of the Forward Rate by no arbitrage}
%\begin{itemize}
%	\item Following Filipovic (2009) we can also directly define the \textcolor{red}{forward rate} by implementing the following:
%	\item at time $t$ sell one $T$-maturity ZCB and buy $\frac{P(t,T)}{P(t,S)}$ ZCBs;
%	\item in $T$ you must pay 1;
%	\item in $S$ you receive $\frac{P(t,T)}{P(t,S)}$, it is equivalent to a forward start investment of 1 in $T$;
%	\item so 
%	\begin{equation*}
	%	1 + (S-T)F(t;T, S)= \frac{P(t,T)}{P(t,S)}
	%	\end{equation*}
%	\item hence we get as before the expression for the forward rate.
%	\item Are there risks in the implementation of this strategy ?
%\end{itemize}	
%\end{frame}
%
%\begin{frame}{FRA: detailed proof}
%We start by
%\begin{equation}
%\mathbb{E}_t^Q[D(t, S)\tau K – D(t,S) \tau L(T,S)]
%\end{equation}
%From which we derive easily
%\begin{align}
%\tau K \mathbb{E}_t^Q[D(t, S)]&-\mathbb{E}_t^Q[D(t, S)\tau L(T,S)]=\\
%& \tau K P(t,S) - \mathbb{E}_t^Q[D(t, S)\tau L(T,S)]
%\end{align}
%as $P(t,S)=\mathbb{E}_t^Q[D(t, S)]$
%
%But remembering the definition given above
%$D(t,S) = D(t,T)D(T,S)$
%\end{frame}
%
%\begin{frame}{}
%Substituting in (1)
%\begin{equation}
%\tau KP(t,S)-\mathbb{E}_t^Q[\tau D(t, T)D(T,S)L(T,S)]
%\end{equation}
%Then using the properties of the expectation operator
%\begin{equation}
%\tau K P(t,S)-\mathbb{E}_t^Q[\mathbb{E}_t^Q[\tau D(t, T)D(T, S)L(T,S)]]
%\end{equation}
%From the last equation we get (using another property of expectation)
%\begin{equation}
%\tau K P(t,S)-\mathbb{E}_t^Q[ \tau D(t, T) L(T,S) \mathbb{E}_t^Q[D(T, S)]]
%\end{equation}
%From which we can write by the definition of $P$
%\begin{equation}
%\tau K P(t,S)-\mathbb{E}_t^Q[ \tau D(t, T) L(T,S) P(T,S)]]
%\end{equation}
%\end{frame}
%
%\begin{frame}{}
%From the last equation I can write by the property of Expectation
%\begin{equation}
%\tau K P(t,S)-\mathbb{E}_t^Q[D(t,T)]+\mathbb{E}_t^Q[D(t, T) \mathbb{E}_t^Q[D(T, S)]]
%\end{equation}
%Perversely splitting 
%\begin{equation}
%\tau K P(t,S)-\mathbb{E}_t^Q[D(t,T)]+\mathbb{E}_t^Q[\mathbb{E}_t^Q[D(t, T) D(T, S)]]
%\end{equation}
%More perversely
%\begin{equation}
%\tau K P(t,S)-\mathbb{E}_t^Q[D(t,T)]+\mathbb{E}_t^Q[\mathbb{E}_t^Q[D(t, S)]
%\end{equation}
%Then applying the lay of iterated expectations we get
%\begin{equation}
%\tau K P(t,S)-\mathbb{E}_t^Q[D(t,T)]+\mathbb{E}_t^Q[D(t, S)]
%\end{equation}
%\end{frame}
%
%\begin{frame}{}
%Finally we get what we are looking for
%\begin{equation}
%\tau K P(t,S)-P(t, T)+ P(t, S)
%\end{equation}
%\end{frame}

%\begin{frame}{Forward rates: shedding light on an underlying relation}
%\begin{itemize}
%\item The starting point is the observation that in a FRA, we must evaluate the future value of the rate $L(T,S)$
%\item Let us define the forward zero coupon bond in valuation date $t$ as
%\begin{equation}
%P(t,T,S):= \mathbb{E}_t^Q[P(T,S)]
%\end{equation}
%\item In order to avoid arbitrage in the calculation of the price $V(t,G(S))$ in $t$ of a generic and certain sum $G(S)$, the following must hold
%\end{itemize}
%\end{frame}
%
%\begin{frame}{}
%\begin{equation}
%V(t,G(T))=P(t,S)G(T)=G(T)P(t,T)P(t,T,S)
%\end{equation}
%Hence we get $P(t,T,S)+\frac{P(t,S)}{P(t,T)}$
%And then we can define $F(t,T,S)$ inverting
%$P(t,T,S)+\frac{P(t,S)}{P(t,T)}:=\frac{1}{1+F(t,T,S)(S-T)}$
%We invert and find the same no arbitrage formula above. The relation
%$P(t,S)=P(t,T)P(t,T,S)$
%Relies on the fact that there is no counterparty risk. Basically the crisis has broken this relationship because the $P(…)$ may be risky: between $t$, $T$ and $S$ many things can happen to one of the participants of the trade.
%\end{frame}

\end{document}
